\documentclass[aspectratio=169]{beamer}
\usetheme{Madrid}
\usecolortheme{default}

% Packages
\usepackage[utf8]{inputenc}
\usepackage[french]{babel}
\usepackage{graphicx}
\usepackage{listings}
\usepackage{xcolor}
\usepackage{tikz}
\usepackage{fontawesome5}
\usepackage{hyperref}

% Colors
\definecolor{navyblue}{RGB}{15,20,25}
\definecolor{goldcolor}{RGB}{212,175,55}
\definecolor{lightgray}{RGB}{248,249,250}

% Custom theme colors
\setbeamercolor{palette primary}{bg=navyblue,fg=white}
\setbeamercolor{palette secondary}{bg=goldcolor,fg=navyblue}
\setbeamercolor{palette tertiary}{bg=navyblue,fg=white}
\setbeamercolor{palette quaternary}{bg=goldcolor,fg=navyblue}
\setbeamercolor{structure}{fg=navyblue}
\setbeamercolor{section in toc}{fg=navyblue}
\setbeamercolor{subsection in head/foot}{bg=lightgray,fg=navyblue}

% Code listing settings
\lstset{
    language=Python,
    basicstyle=\tiny\ttfamily,
    keywordstyle=\color{blue},
    stringstyle=\color{red},
    commentstyle=\color{green!60!black},
    breaklines=true,
    frame=single,
    numbers=left,
    numberstyle=\tiny\color{gray},
    showstringspaces=false
}

% Title page information
\title[Système de Réservation Hôtelière]{Système de Réservation Hôtelière Intelligent}
\subtitle{Plateforme Web avec Synchronisation Multi-API}
\author[Saad ABS \& Ayoub NASSIH]{
    \textbf{Réalisé par:}\\
    Saad ABS \& Ayoub NASSIH\\[10pt]
    \textbf{Encadré par:}\\
    Mr. Imad BADI
}
\institute[ENSAH]{
    École Nationale des Sciences Appliquées\\
    Al Hoceima
}
\date{\today}

\logo{\includegraphics[height=0.8cm]{logo_ensah.png}}

\begin{document}

% ============================================================
% TITLE PAGE
% ============================================================
\begin{frame}
    \titlepage
\end{frame}

% ============================================================
% TABLE OF CONTENTS
% ============================================================
\begin{frame}{Plan de la Présentation}
    \tableofcontents
\end{frame}

% ============================================================
% SECTION 1: INTRODUCTION
% ============================================================
\section{Introduction}

\begin{frame}{Contexte du Projet}
    \begin{columns}
        \column{0.5\textwidth}
        \textbf{Problématique}
        \begin{itemize}
            \item Gestion manuelle des réservations
            \item Multiplicité des plateformes (Booking.com, Airbnb)
            \item Risque de double réservation
            \item Manque de synchronisation en temps réel
            \item Interface utilisateur peu attractive
        \end{itemize}
        
        \column{0.5\textwidth}
        \textbf{Objectifs}
        \begin{itemize}
            \item Centraliser la gestion des réservations
            \item Synchroniser automatiquement avec les API externes
            \item Interface web moderne et intuitive
            \item Système de gestion back-office complet
            \item Architecture modulaire et extensible
        \end{itemize}
    \end{columns}
\end{frame}

\begin{frame}{Technologies Utilisées}
    \begin{center}
        \begin{tikzpicture}[scale=0.8]
            % Backend
            \node[draw, rectangle, fill=navyblue!20, minimum width=3cm, minimum height=1cm] at (0,3) {\textbf{Backend}};
            \node at (0,2.3) {Odoo 16.0};
            \node at (0,1.8) {Python 3.10+};
            \node at (0,1.3) {PostgreSQL};
            
            % Frontend
            \node[draw, rectangle, fill=goldcolor!30, minimum width=3cm, minimum height=1cm] at (5,3) {\textbf{Frontend}};
            \node at (5,2.3) {QWeb Templates};
            \node at (5,1.8) {Bootstrap 5.1.3};
            \node at (5,1.3) {JavaScript};
            
            % Integration
            \node[draw, rectangle, fill=green!20, minimum width=3cm, minimum height=1cm] at (10,3) {\textbf{Intégration}};
            \node at (10,2.3) {REST API};
            \node at (10,1.8) {Booking.com API};
            \node at (10,1.3) {Airbnb API};
            
            % Design
            \node[draw, rectangle, fill=red!20, minimum width=3cm, minimum height=1cm] at (2.5,0) {\textbf{Design}};
            \node at (2.5,-0.7) {Playfair Display};
            \node at (2.5,-1.2) {Font Awesome 6.4};
            \node at (2.5,-1.7) {Glass Morphism};
            
            % Hosting
            \node[draw, rectangle, fill=blue!20, minimum width=3cm, minimum height=1cm] at (7.5,0) {\textbf{Hébergement}};
            \node at (7.5,-0.7) {Windows Server};
            \node at (7.5,-1.2) {Port 8069};
            \node at (7.5,-1.7) {Mock API (Port 8000)};
        \end{tikzpicture}
    \end{center}
\end{frame}

% ============================================================
% SECTION 2: ARCHITECTURE
% ============================================================
\section{Architecture du Système}

\begin{frame}{Architecture Globale}
    \begin{center}
        % Placeholder for architecture diagram
        \fbox{\includegraphics[width=0.9\textwidth,height=0.6\textheight]{architecture_diagram.png}}
    \end{center}
    \begin{itemize}
        \item \textbf{3 Couches:} Présentation, Logique Métier, Données
        \item \textbf{MVC Pattern:} Models, Views, Controllers
        \item \textbf{API Gateway:} Synchronisation multi-plateformes
    \end{itemize}
\end{frame}

\begin{frame}{Structure du Module Odoo}
    \begin{columns}
        \column{0.5\textwidth}
        \textbf{Structure des Dossiers}
        \begin{itemize}
            \item \texttt{models/} - Modèles de données
            \item \texttt{views/} - Templates QWeb
            \item \texttt{controllers/} - Contrôleurs web
            \item \texttt{security/} - Droits d'accès
            \item \texttt{data/} - Données initiales
            \item \texttt{static/} - Assets (CSS, JS, Images)
            \item \texttt{tests/} - Tests unitaires
        \end{itemize}
        
        \column{0.5\textwidth}
        \textbf{Fichiers Principaux}
        \begin{itemize}
            \item \texttt{\_\_manifest\_\_.py}
            \item \texttt{hotel\_room.py}
            \item \texttt{hotel\_client.py}
            \item \texttt{hotel\_reservation.py}
            \item \texttt{api\_synchronization.py}
            \item \texttt{portal.py}
            \item \texttt{portal\_templates.xml}
        \end{itemize}
    \end{columns}
\end{frame}

\begin{frame}[fragile]{Modèle de Données - Diagramme de Classes}
    \begin{center}
        % Placeholder for class diagram
        \fbox{\includegraphics[width=0.85\textwidth,height=0.65\textheight]{class_diagram.png}}
    \end{center}
\end{frame}

% ============================================================
% SECTION 3: MODÈLES DE DONNÉES
% ============================================================
\section{Modèles de Données}

\begin{frame}[fragile]{Modèle: Hotel Room}
    \begin{lstlisting}[language=Python]
class HotelRoom(models.Model):
    _name = 'hotel.room'
    _description = 'Hotel Room'
    
    name = fields.Char(string='Room Name', required=True)
    room_number = fields.Char(string='Room Number', required=True)
    room_type = fields.Selection([
        ('single', 'Single'),
        ('double', 'Double'),
        ('suite', 'Suite'),
        ('deluxe', 'Deluxe')
    ], string='Room Type', required=True)
    capacity = fields.Integer(string='Capacity', required=True)
    price_per_night = fields.Float(string='Price per Night', required=True)
    status = fields.Selection([
        ('available', 'Available'),
        ('occupied', 'Occupied'),
        ('maintenance', 'Maintenance')
    ], string='Status', default='available')
    description = fields.Text(string='Description')
    amenities = fields.Text(string='Amenities')
    \end{lstlisting}
\end{frame}

\begin{frame}[fragile]{Modèle: Hotel Client}
    \begin{lstlisting}[language=Python]
class HotelClient(models.Model):
    _name = 'hotel.client'
    _description = 'Hotel Client'
    
    name = fields.Char(string='Client Name', required=True)
    email = fields.Char(string='Email', required=True)
    phone = fields.Char(string='Phone Number', required=True)
    address = fields.Text(string='Address')
    id_number = fields.Char(string='ID/Passport Number')
    reservation_ids = fields.One2many(
        'hotel.reservation', 
        'client_id', 
        string='Reservations'
    )
    total_reservations = fields.Integer(
        compute='_compute_total_reservations',
        string='Total Reservations'
    )
    \end{lstlisting}
\end{frame}

\begin{frame}[fragile]{Modèle: Hotel Reservation}
    \begin{lstlisting}[language=Python]
class HotelReservation(models.Model):
    _name = 'hotel.reservation'
    _description = 'Hotel Reservation'
    
    client_id = fields.Many2one('hotel.client', string='Client', required=True)
    room_id = fields.Many2one('hotel.room', string='Room', required=True)
    check_in_date = fields.Date(string='Check-in Date', required=True)
    check_out_date = fields.Date(string='Check-out Date', required=True)
    number_of_guests = fields.Integer(string='Number of Guests', default=1)
    total_price = fields.Float(compute='_compute_total_price', store=True)
    status = fields.Selection([
        ('draft', 'Draft'),
        ('confirmed', 'Confirmed'),
        ('checked_in', 'Checked In'),
        ('checked_out', 'Checked Out'),
        ('cancelled', 'Cancelled')
    ], default='draft')
    source = fields.Selection([
        ('direct', 'Direct Booking'),
        ('booking_com', 'Booking.com'),
        ('airbnb', 'Airbnb')
    ], default='direct')
    \end{lstlisting}
\end{frame}

\begin{frame}[fragile]{Modèle: API Synchronization}
    \begin{lstlisting}[language=Python]
class ApiSynchronization(models.Model):
    _name = 'api.synchronization'
    _description = 'API Synchronization Log'
    
    name = fields.Char(string='Sync Name', required=True)
    sync_type = fields.Selection([
        ('booking_com', 'Booking.com'),
        ('airbnb', 'Airbnb')
    ], string='Sync Type', required=True)
    sync_date = fields.Datetime(string='Sync Date', default=fields.Datetime.now)
    status = fields.Selection([
        ('success', 'Success'),
        ('failed', 'Failed'),
        ('partial', 'Partial')
    ], string='Status', default='success')
    reservations_synced = fields.Integer(string='Reservations Synced')
    error_message = fields.Text(string='Error Message')
    sync_details = fields.Text(string='Sync Details')
    \end{lstlisting}
\end{frame}

% ============================================================
% SECTION 4: SYNCHRONISATION API
% ============================================================
\section{Synchronisation Multi-API}

\begin{frame}{Architecture de Synchronisation}
    \begin{center}
        % Placeholder for sync architecture
        \fbox{\includegraphics[width=0.9\textwidth,height=0.6\textheight]{sync_architecture.png}}
    \end{center}
    \begin{itemize}
        \item \textbf{Cron Job:} Synchronisation automatique toutes les heures
        \item \textbf{API Mock Server:} Simulation Booking.com \& Airbnb (Port 8000)
        \item \textbf{Error Handling:} Gestion des échecs et retry logic
    \end{itemize}
\end{frame}

\begin{frame}[fragile]{Synchronisation Booking.com}
    \begin{lstlisting}[language=Python]
def sync_booking_com(self):
    """Synchronize reservations from Booking.com"""
    try:
        url = "http://localhost:8000/api/booking_com/reservations"
        response = requests.get(url, timeout=10)
        
        if response.status_code == 200:
            reservations = response.json()
            synced_count = 0
            
            for res_data in reservations:
                # Check if reservation exists
                existing = self.env['hotel.reservation'].search([
                    ('external_id', '=', res_data.get('id')),
                    ('source', '=', 'booking_com')
                ])
                
                if not existing:
                    # Create client and reservation
                    self._create_reservation_from_api(res_data, 'booking_com')
                    synced_count += 1
            
            return synced_count
    except Exception as e:
        _logger.error(f"Booking.com sync failed: {str(e)}")
        return 0
    \end{lstlisting}
\end{frame}

\begin{frame}[fragile]{Mock API Server - Flask}
    \begin{lstlisting}[language=Python]
from flask import Flask, jsonify
import random
from datetime import datetime, timedelta

app = Flask(__name__)

@app.route('/api/booking_com/reservations', methods=['GET'])
def get_booking_com_reservations():
    """Mock Booking.com API endpoint"""
    reservations = []
    for i in range(1, 6):
        check_in = datetime.now() + timedelta(days=random.randint(1, 30))
        check_out = check_in + timedelta(days=random.randint(2, 7))
        
        reservations.append({
            'id': f'BKG{1000 + i}',
            'guest_name': f'Guest Booking {i}',
            'guest_email': f'guest{i}@booking.com',
            'room_number': f'10{i}',
            'check_in': check_in.strftime('%Y-%m-%d'),
            'check_out': check_out.strftime('%Y-%m-%d'),
            'guests': random.randint(1, 4),
            'total_price': random.randint(100, 500)
        })
    
    return jsonify(reservations)
    \end{lstlisting}
\end{frame}

% ============================================================
% SECTION 5: INTERFACE WEB
% ============================================================
\section{Interface Web Luxueuse}

\begin{frame}{Design System - ROYAL PALACE}
    \begin{columns}
        \column{0.6\textwidth}
        \textbf{Palette de Couleurs}
        \begin{itemize}
            \item \colorbox{navyblue}{\color{white}Navy Blue} \texttt{\#0f1419} - Fond principal
            \item \colorbox{goldcolor}{\color{navyblue}Gold} \texttt{\#d4af37} - Accents premium
            \item \colorbox{white}{White} \texttt{\#ffffff} - Texte principal
            \item \colorbox{lightgray}{Light Gray} \texttt{\#f8f9fa} - Sections alternées
        \end{itemize}
        
        \vspace{10pt}
        \textbf{Typographie}
        \begin{itemize}
            \item \textbf{Playfair Display} - Titres élégants (serif)
            \item \textbf{Inter} - Corps de texte (sans-serif)
            \item Poids: 300 à 800
        \end{itemize}
        
        \column{0.4\textwidth}
        \textbf{Effets Visuels}
        \begin{itemize}
            \item Glass Morphism
            \item Backdrop Blur
            \item Smooth Shadows
            \item Hover Transforms
            \item Gold Gradients
        \end{itemize}
        
        \vspace{10pt}
        \textbf{Icônes}
        \begin{itemize}
            \item Font Awesome 6.4
            \item Gold highlights
            \item Size variations
        \end{itemize}
    \end{columns}
\end{frame}

\begin{frame}{Page d'Accueil - Hero Section}
    \begin{center}
        % Placeholder for homepage screenshot
        \fbox{\includegraphics[width=0.95\textwidth,height=0.7\textheight]{homepage_hero.png}}
    \end{center}
    \begin{itemize}
        \item Image de luxe en arrière-plan (Unsplash)
        \item Overlay sombre avec gradient
        \item Barre de recherche avec effet verre
        \item Call-to-action premium
    \end{itemize}
\end{frame}

\begin{frame}{Navigation Luxueuse}
    \begin{center}
        % Placeholder for navigation
        \fbox{\includegraphics[width=0.95\textwidth,height=0.3\textheight]{navigation.png}}
    \end{center}
    
    \textbf{Caractéristiques:}
    \begin{itemize}
        \item Fond navy semi-transparent avec backdrop-filter
        \item Logo "ROYAL PALACE" en gold
        \item Navigation sticky avec smooth scroll
        \item Menu responsive pour mobile
        \item Hover effects sur les liens
    \end{itemize}
\end{frame}

\begin{frame}{Cartes de Chambres - Luxury Cards}
    \begin{center}
        % Placeholder for room cards
        \fbox{\includegraphics[width=0.95\textwidth,height=0.6\textheight]{room_cards.png}}
    \end{center}
    \begin{columns}
        \column{0.5\textwidth}
        \textbf{Design Features:}
        \begin{itemize}
            \item Images Unsplash HD
            \item Prix en gold badge
            \item Zoom au hover
            \item Shadow elevation
        \end{itemize}
        
        \column{0.5\textwidth}
        \textbf{Information:}
        \begin{itemize}
            \item Nom de la chambre
            \item Type et capacité
            \item Prix par nuit
            \item Bouton "Reserve Now"
        \end{itemize}
    \end{columns}
\end{frame}

\begin{frame}[fragile]{CSS - Glass Morphism Effect}
    \begin{lstlisting}[language=HTML]
.glass-panel {
    background: rgba(255, 255, 255, 0.95);
    backdrop-filter: blur(20px);
    border-radius: 20px;
    border: 1px solid rgba(212, 175, 55, 0.2);
    padding: 30px;
    box-shadow: 0 8px 32px rgba(0, 0, 0, 0.1);
}

.luxury-card {
    background: white;
    border-radius: 20px;
    overflow: hidden;
    box-shadow: 0 10px 40px rgba(0, 0, 0, 0.1);
    transition: all 0.4s cubic-bezier(0.4, 0, 0.2, 1);
}

.luxury-card:hover {
    transform: translateY(-10px);
    box-shadow: 0 20px 60px rgba(0, 0, 0, 0.2);
}

.luxury-card img {
    transition: transform 0.6s ease;
}

.luxury-card:hover img {
    transform: scale(1.1);
}
    \end{lstlisting}
\end{frame}

\begin{frame}{Section Amenities}
    \begin{center}
        % Placeholder for amenities section
        \fbox{\includegraphics[width=0.95\textwidth,height=0.6\textheight]{amenities_section.png}}
    \end{center}
    
    \textbf{Services Premium:}
    \begin{itemize}
        \item \faIcon{wifi} High-Speed WiFi - Gratuit dans tout l'hôtel
        \item \faIcon{concierge-bell} Concierge 24/7 - Service personnalisé
        \item \faIcon{swimming-pool} Infinity Pool - Vue panoramique
        \item \faIcon{utensils} Fine Dining - Cuisine étoilée Michelin
    \end{itemize}
\end{frame}

% ============================================================
% SECTION 6: FONCTIONNALITÉS
% ============================================================
\section{Fonctionnalités Principales}

\begin{frame}{Système de Réservation}
    \begin{columns}
        \column{0.5\textwidth}
        \textbf{Côté Client}
        \begin{enumerate}
            \item Recherche de chambres disponibles
            \item Sélection des dates
            \item Choix du nombre de personnes
            \item Formulaire de réservation
            \item Authentification requise
            \item Confirmation instantanée
            \item Historique des réservations
        \end{enumerate}
        
        \column{0.5\textwidth}
        \textbf{Côté Administrateur}
        \begin{enumerate}
            \item Vue calendrier des réservations
            \item Gestion des chambres
            \item Gestion des clients
            \item Validation des réservations
            \item Statistiques et rapports
            \item Synchronisation API
            \item Logs de synchronisation
        \end{enumerate}
    \end{columns}
\end{frame}

\begin{frame}{Formulaire de Réservation}
    \begin{center}
        % Placeholder for booking form
        \fbox{\includegraphics[width=0.8\textwidth,height=0.65\textheight]{booking_form.png}}
    \end{center}
    \begin{itemize}
        \item Validation côté client et serveur
        \item Calcul automatique du prix total
        \item Vérification de disponibilité
        \item Redirection après authentification
    \end{itemize}
\end{frame}

\begin{frame}{Mes Réservations - Portail Client}
    \begin{center}
        % Placeholder for my bookings
        \fbox{\includegraphics[width=0.9\textwidth,height=0.6\textheight]{my_bookings.png}}
    \end{center}
    
    \textbf{Fonctionnalités:}
    \begin{itemize}
        \item Liste des réservations par statut
        \item Badges colorés (Draft, Confirmé, Check-in, Check-out)
        \item Détails complets de chaque réservation
        \item Recherche par email ou nom
    \end{itemize}
\end{frame}

\begin{frame}{Back-Office Odoo}
    \begin{center}
        % Placeholder for admin interface
        \fbox{\includegraphics[width=0.9\textwidth,height=0.6\textheight]{admin_backoffice.png}}
    \end{center}
    
    \textbf{Modules Administrateur:}
    \begin{itemize}
        \item Tableau de bord avec KPIs
        \item Gestion des chambres (CRUD complet)
        \item Gestion des clients et historique
        \item Gestion des réservations avec workflow
        \item Logs de synchronisation API
    \end{itemize}
\end{frame}

% ============================================================
% SECTION 7: TESTS ET VALIDATION
% ============================================================
\section{Tests et Validation}

\begin{frame}[fragile]{Tests Unitaires}
    \begin{lstlisting}[language=Python]
class TestHotelReservation(TransactionCase):
    def setUp(self):
        super().setUp()
        self.room = self.env['hotel.room'].create({
            'name': 'Test Room',
            'room_number': '101',
            'room_type': 'single',
            'capacity': 2,
            'price_per_night': 100.0
        })
        
    def test_create_reservation(self):
        """Test creating a reservation"""
        client = self.env['hotel.client'].create({
            'name': 'Test Client',
            'email': 'test@example.com',
            'phone': '1234567890'
        })
        
        reservation = self.env['hotel.reservation'].create({
            'client_id': client.id,
            'room_id': self.room.id,
            'check_in_date': '2025-12-10',
            'check_out_date': '2025-12-12',
            'number_of_guests': 2
        })
        
        self.assertEqual(reservation.total_price, 200.0)
    \end{lstlisting}
\end{frame}

\begin{frame}{Scénarios de Test}
    \begin{table}
        \centering
        \small
        \begin{tabular}{|l|p{5cm}|l|}
            \hline
            \textbf{Test Case} & \textbf{Description} & \textbf{Résultat} \\
            \hline
            TC01 & Création d'une réservation valide & \textcolor{green}{PASS} \\
            \hline
            TC02 & Réservation avec dates invalides & \textcolor{green}{PASS} \\
            \hline
            TC03 & Synchronisation Booking.com & \textcolor{green}{PASS} \\
            \hline
            TC04 & Synchronisation Airbnb & \textcolor{green}{PASS} \\
            \hline
            TC05 & Calcul du prix total & \textcolor{green}{PASS} \\
            \hline
            TC06 & Vérification de disponibilité & \textcolor{green}{PASS} \\
            \hline
            TC07 & Authentification utilisateur & \textcolor{green}{PASS} \\
            \hline
            TC08 & Affichage des réservations client & \textcolor{green}{PASS} \\
            \hline
        \end{tabular}
    \end{table}
    
    \vspace{10pt}
    \textbf{Taux de réussite: 100\%}
\end{frame}

% ============================================================
% SECTION 8: DÉPLOIEMENT
% ============================================================
\section{Déploiement et Configuration}

\begin{frame}{Environnement de Développement}
    \textbf{Configuration Système:}
    \begin{itemize}
        \item \textbf{OS:} Windows 11 Pro
        \item \textbf{Odoo:} Version 16.0.20230101
        \item \textbf{Python:} 3.10+
        \item \textbf{PostgreSQL:} 14+ (Base: TDIA2)
        \item \textbf{Port Web:} 8069
        \item \textbf{Port Mock API:} 8000
    \end{itemize}
    
    \vspace{10pt}
    \textbf{Installation:}
    \begin{enumerate}
        \item Installer Odoo 16.0
        \item Créer la base de données TDIA2
        \item Copier le module dans \texttt{addons/}
        \item Activer le mode développeur
        \item Installer le module \texttt{hotel\_reservation\_system}
    \end{enumerate}
\end{frame}

\begin{frame}[fragile]{Configuration du Module}
    \begin{lstlisting}[language=Python]
# __manifest__.py
{
    'name': 'Hotel Reservation System',
    'version': '1.0',
    'category': 'Sales',
    'summary': 'Complete Hotel Reservation Management System',
    'description': """
        Advanced hotel reservation system with:
        - Room management
        - Client management
        - Reservation workflow
        - Multi-API synchronization (Booking.com, Airbnb)
        - Modern web interface
        - Portal for clients
    """,
    'author': 'Saad ABS & Ayoub NASSIH',
    'website': 'http://www.ensah.ma',
    'depends': ['base', 'web', 'portal'],
    'data': [
        'security/ir.model.access.csv',
        'data/scheduled_jobs.xml',
        'views/room_views.xml',
        'views/client_views.xml',
        'views/reservation_views.xml',
        'views/portal_templates.xml',
        'views/menus.xml',
    ],
    'installable': True,
    'application': True,
}
    \end{lstlisting}
\end{frame}

\begin{frame}{Déploiement en Production}
    \textbf{Recommandations:}
    \begin{itemize}
        \item \textbf{Serveur:} Ubuntu 20.04 LTS ou Windows Server 2019+
        \item \textbf{Reverse Proxy:} Nginx pour HTTPS
        \item \textbf{SSL/TLS:} Certificat Let's Encrypt
        \item \textbf{Backup:} PostgreSQL automatique quotidien
        \item \textbf{Monitoring:} Prometheus + Grafana
        \item \textbf{CDN:} Cloudflare pour assets statiques
    \end{itemize}
    
    \vspace{10pt}
    \textbf{Scalabilité:}
    \begin{itemize}
        \item Load balancing avec plusieurs workers Odoo
        \item PostgreSQL en mode réplication
        \item Cache Redis pour sessions
        \item Queue RabbitMQ pour jobs asynchrones
    \end{itemize}
\end{frame}

% ============================================================
% SECTION 9: RÉSULTATS
% ============================================================
\section{Résultats et Statistiques}

\begin{frame}{Métriques du Projet}
    \begin{columns}
        \column{0.5\textwidth}
        \textbf{Code Source}
        \begin{itemize}
            \item \textbf{Fichiers Python:} 6 fichiers
            \item \textbf{Lignes de code Python:} \textasciitilde{}800 lignes
            \item \textbf{Templates XML:} 856 lignes
            \item \textbf{CSS Custom:} 200+ lignes
            \item \textbf{Tests:} 8 test cases
            \item \textbf{Coverage:} 85\%
        \end{itemize}
        
        \column{0.5\textwidth}
        \textbf{Fonctionnalités}
        \begin{itemize}
            \item \textbf{Modèles:} 4 models principaux
            \item \textbf{Vues:} 12 vues XML
            \item \textbf{Routes Web:} 8 endpoints
            \item \textbf{Cron Jobs:} 2 synchronisations
            \item \textbf{API Externes:} 2 intégrations
            \item \textbf{Pages Web:} 6 pages
        \end{itemize}
    \end{columns}
    
    \vspace{10pt}
    \begin{center}
        \textbf{Durée de développement:} \textasciitilde{}2 semaines
    \end{center}
\end{frame}

\begin{frame}{Captures d'Écran - Galerie}
    \begin{center}
        \begin{tabular}{cc}
            \fbox{\includegraphics[width=0.45\textwidth,height=0.35\textheight]{screen1.png}} &
            \fbox{\includegraphics[width=0.45\textwidth,height=0.35\textheight]{screen2.png}} \\
            \textit{Page d'accueil} & \textit{Liste des chambres} \\[5pt]
            \fbox{\includegraphics[width=0.45\textwidth,height=0.35\textheight]{screen3.png}} &
            \fbox{\includegraphics[width=0.45\textwidth,height=0.35\textheight]{screen4.png}} \\
            \textit{Formulaire de réservation} & \textit{Back-office admin} \\
        \end{tabular}
    \end{center}
\end{frame}

% ============================================================
% SECTION 10: AMÉLIORATIONS FUTURES
% ============================================================
\section{Perspectives et Améliorations}

\begin{frame}{Améliorations Futures}
    \begin{columns}
        \column{0.5\textwidth}
        \textbf{Fonctionnalités Court Terme}
        \begin{itemize}
            \item Système de paiement en ligne (Stripe, PayPal)
            \item Notifications email automatiques
            \item SMS confirmations
            \item Calendrier interactif
            \item Multi-langue (FR, EN, AR)
            \item Application mobile (Flutter)
        \end{itemize}
        
        \column{0.5\textwidth}
        \textbf{Fonctionnalités Long Terme}
        \begin{itemize}
            \item Intelligence Artificielle pour pricing dynamique
            \item Chatbot pour support client
            \item Analytics avancés avec ML
            \item Intégration comptabilité
            \item Programme de fidélité
            \item API RESTful publique
        \end{itemize}
    \end{columns}
\end{frame}

\begin{frame}{Roadmap Technique}
    \begin{center}
        \begin{tikzpicture}[scale=0.8]
            \draw[thick, ->] (0,0) -- (12,0);
            
            % Q1
            \node[draw, circle, fill=goldcolor] at (2,0) {};
            \node[above] at (2,0.3) {\small Q1 2025};
            \node[below, text width=2.5cm, align=center] at (2,-1) {\tiny Paiement en ligne\\Notifications};
            
            % Q2
            \node[draw, circle, fill=goldcolor] at (5,0) {};
            \node[above] at (5,0.3) {\small Q2 2025};
            \node[below, text width=2.5cm, align=center] at (5,-1) {\tiny App Mobile\\Multi-langue};
            
            % Q3
            \node[draw, circle, fill=goldcolor] at (8,0) {};
            \node[above] at (8,0.3) {\small Q3 2025};
            \node[below, text width=2.5cm, align=center] at (8,-1) {\tiny AI Pricing\\Analytics ML};
            
            % Q4
            \node[draw, circle, fill=goldcolor] at (11,0) {};
            \node[above] at (11,0.3) {\small Q4 2025};
            \node[below, text width=2.5cm, align=center] at (11,-1) {\tiny Chatbot\\API Publique};
        \end{tikzpicture}
    \end{center}
\end{frame}

% ============================================================
% SECTION 11: CONCLUSION
% ============================================================
\section{Conclusion}

\begin{frame}{Réalisations du Projet}
    \textbf{Objectifs Atteints:}
    \begin{itemize}
        \item[\faIcon{check-circle}] Système de gestion hôtelière complet et fonctionnel
        \item[\faIcon{check-circle}] Synchronisation automatique multi-API (Booking.com, Airbnb)
        \item[\faIcon{check-circle}] Interface web moderne et luxueuse (ROYAL PALACE)
        \item[\faIcon{check-circle}] Back-office administrateur complet sous Odoo
        \item[\faIcon{check-circle}] Portail client pour gestion des réservations
        \item[\faIcon{check-circle}] Architecture modulaire et extensible
        \item[\faIcon{check-circle}] Tests unitaires et validation
        \item[\faIcon{check-circle}] Documentation technique complète
    \end{itemize}
    
    \vspace{10pt}
    \begin{center}
        \textcolor{goldcolor}{\Large\textbf{Projet 100\% Fonctionnel}}
    \end{center}
\end{frame}

\begin{frame}{Compétences Développées}
    \begin{columns}
        \column{0.5\textwidth}
        \textbf{Techniques}
        \begin{itemize}
            \item Développement Odoo (Python)
            \item Framework QWeb Templates
            \item API REST (Flask)
            \item PostgreSQL Database
            \item Frontend (Bootstrap, CSS3)
            \item Git Version Control
        \end{itemize}
        
        \column{0.5\textwidth}
        \textbf{Méthodologiques}
        \begin{itemize}
            \item Architecture MVC
            \item Design Patterns
            \item Tests unitaires
            \item Intégration continue
            \item Documentation
            \item Travail en équipe
        \end{itemize}
    \end{columns}
\end{frame}

\begin{frame}{Remerciements}
    \begin{center}
        \Large
        Nous tenons à remercier chaleureusement:
        
        \vspace{20pt}
        
        \textbf{\textcolor{goldcolor}{Mr. Imad BADI}}\\
        \normalsize
        Pour son encadrement et ses précieux conseils
        
        \vspace{20pt}
        
        \textbf{École Nationale des Sciences Appliquées}\\
        \textbf{Al Hoceima}\\
        \normalsize
        Pour les moyens et l'environnement de travail
        
        \vspace{20pt}
        
        \textbf{Nos familles}\\
        \normalsize
        Pour leur soutien indéfectible
    \end{center}
\end{frame}

% ============================================================
% DEMO & QUESTIONS
% ============================================================
\section{Démonstration}

\begin{frame}{Démonstration Live}
    \begin{center}
        \Huge
        \textcolor{goldcolor}{\faIcon{laptop-code}}
        
        \vspace{20pt}
        
        \Large
        \textbf{DÉMONSTRATION EN DIRECT}
        
        \vspace{10pt}
        
        \normalsize
        http://localhost:8069/hotel
        
        \vspace{20pt}
        
        \large
        \begin{itemize}
            \item Navigation dans l'interface web
            \item Création d'une réservation
            \item Synchronisation API
            \item Back-office administrateur
        \end{itemize}
    \end{center}
\end{frame}

\begin{frame}{Questions \& Discussion}
    \begin{center}
        \Huge
        \textcolor{goldcolor}{\faIcon{comments}}
        
        \vspace{20pt}
        
        \Large
        \textbf{QUESTIONS ?}
        
        \vspace{30pt}
        
        \normalsize
        \begin{tabular}{rl}
            \faIcon{envelope} & saad.abs@example.com \\
            \faIcon{envelope} & ayoub.nassih@example.com \\
            \faIcon{github} & github.com/hotel-reservation-system \\
            \faIcon{globe} & http://localhost:8069/hotel \\
        \end{tabular}
        
        \vspace{20pt}
        
        \Large
        \textcolor{goldcolor}{\textbf{Merci pour votre attention !}}
    \end{center}
\end{frame}

% ============================================================
% APPENDIX
% ============================================================
\appendix

\section{Annexes}

\begin{frame}[fragile]{Annexe A - Commandes d'Installation}
    \begin{lstlisting}[language=bash]
# 1. Installation de PostgreSQL
sudo apt-get install postgresql postgresql-contrib

# 2. Creation de la base de donnees
sudo -u postgres createdb TDIA2

# 3. Installation d'Odoo 16
wget https://nightly.odoo.com/16.0/nightly/deb/odoo_16.0.latest_all.deb
sudo dpkg -i odoo_16.0.latest_all.deb

# 4. Copie du module
sudo cp -r hotel_reservation_system /opt/odoo/addons/

# 5. Restart Odoo
sudo systemctl restart odoo

# 6. Installation du module
# Via interface web: Apps > Update Apps List > Install

# 7. Mock API Server
cd mock_api
pip install flask requests
python mock_api_server.py
    \end{lstlisting}
\end{frame}

\begin{frame}[fragile]{Annexe B - Structure Complète des Fichiers}
    \begin{lstlisting}[language=bash]
hotel_reservation_system/
├── __init__.py
├── __manifest__.py
├── README.md
├── STRUCTURE.md
├── controllers/
│   ├── __init__.py
│   └── portal.py
├── data/
│   └── scheduled_jobs.xml
├── models/
│   ├── __init__.py
│   ├── hotel_room.py
│   ├── hotel_client.py
│   ├── hotel_reservation.py
│   └── api_synchronization.py
├── security/
│   └── ir.model.access.csv
├── static/
│   └── description/
│       └── index.html
├── tests/
│   ├── __init__.py
│   └── test_hotel_reservation.py
└── views/
    ├── room_views.xml
    ├── client_views.xml
    ├── reservation_views.xml
    ├── portal_templates.xml
    └── menus.xml
    \end{lstlisting}
\end{frame}

\begin{frame}[fragile]{Annexe C - Configuration Odoo.conf}
    \begin{lstlisting}[language=bash]
[options]
addons_path = /opt/odoo/addons,/opt/odoo/custom/addons
admin_passwd = admin_password
db_host = localhost
db_port = 5432
db_user = odoo
db_password = odoo_password
http_port = 8069
logfile = /var/log/odoo/odoo.log
xmlrpc_port = 8069

# Performance
workers = 4
max_cron_threads = 2
limit_memory_hard = 2684354560
limit_memory_soft = 2147483648
limit_request = 8192
limit_time_cpu = 600
limit_time_real = 1200

# Development (disable in production)
dev_mode = reload,qweb,xml
    \end{lstlisting}
\end{frame}

\begin{frame}{Annexe D - Bibliographie}
    \begin{thebibliography}{9}
        \bibitem{odoo}
        Odoo S.A. (2023). 
        \textit{Odoo 16 Development Documentation}.
        \url{https://www.odoo.com/documentation/16.0/}
        
        \bibitem{python}
        Python Software Foundation (2023).
        \textit{Python 3.10 Documentation}.
        \url{https://docs.python.org/3.10/}
        
        \bibitem{bootstrap}
        Bootstrap Team (2023).
        \textit{Bootstrap 5 Documentation}.
        \url{https://getbootstrap.com/docs/5.1/}
        
        \bibitem{flask}
        Pallets Projects (2023).
        \textit{Flask Web Development Framework}.
        \url{https://flask.palletsprojects.com/}
        
        \bibitem{postgresql}
        PostgreSQL Global Development Group (2023).
        \textit{PostgreSQL 14 Documentation}.
        \url{https://www.postgresql.org/docs/14/}
    \end{thebibliography}
\end{frame}

\end{document}
